
La fuente será considerada como desecho radiactivo si su recuperación resulta inconveniente o, bien, si presenta deformidades o fugas que impidan su uso.

Los procedimientos que hemos descrito no constituyen la única manera de resolver un problema y no deben intentarse a menos que el personal esté familiarizado con ellos y los haya practicado en simulacros. La misma recomendación es válida para cualquier procedimiento que intenten emprenderse. Si por algún motivo el técnico no estuviera familiarizado con los procedimientos que debe seguir o no se sintiera lo suficientemente seguro para ejecutarlos, deberá mantener la zona restringida debidamente señalada, acordonada, despejada y vigilada, hasta que llegue a la zona personal capacitado para realizar las operaciones.

La pérdida de fuentes radioactivas será considerada siempre como emergencia radiológica, dado que puede haber vidas involucradas de por medio. Las acciones inmediatas que se deben seguir siempre son:
\begin{itemize}
    \item Notificar al encargado de la seguridad radiológica.
    \item Notificar a la C.N.S.N.S.
    \item Notificar a las personas del área sobre el hecho.
    \item Iniciar la búsqueda de la fuente con ayuda del equipo detector.
\end{itemize}

Las acciones posteriores a seguir, dependerán de cada caso en particular.

Para concluir, diremos que cualquier ahorro en protección radiológica aumenta considerablemente la probabilidad de que se presente un accidente.
